\chapter{Huit problèmes de synthèse}
\section*{Problème S1 -- Une boîte et son coût}

Une feuille carrée de côté $12$ cm est découpée aux quatre coins en carrés
de côté $x$, puis repliée en boîte sans couvercle.
\begin{enumerate}
\item Définir le domaine physique et le volume $V(x)$.
\item Déterminer les dimensions donnant le volume maximal.
\item La surface conservée est facturée à un prix proportionnel à son aire.
Calculer cette aire au maximum et distinguer cette optimisation de celle du coût seul.
\item Une machine choisit $x=2+h$. Donner exactement $V(2+h)-V(2)$,
puis borner la perte de volume pour $|h|\le0{,}1$.
\end{enumerate}

\section*{Problème S2 -- Stock : temps discret et temps continu}

Un stock idéal évolue selon $u_{n+1}=0{,}8u_n+30$, $u_0=50$.
\begin{enumerate}
\item Trouver son équilibre, sa formule explicite et ses variations.
\item Déterminer le premier rang où $u_n\ge140$, par calcul ou programme borné.
\item Comparer avec le modèle continu $y'=-0{,}2y+30$, $y(0)=50$.
\item Trouver le coefficient $k>0$ pour que la solution de
$z'=-k(z-150)$, $z(0)=50$, reproduise exactement $u_n$ aux temps entiers.
\end{enumerate}

\section*{Problème S3 -- Projection et cercle}

On donne $M=(4,3)$ et la droite $d:x+2y-5=0$.
\begin{enumerate}
\item Calculer le projeté $H$ de $M$ et la distance à $d$.
\item Déterminer les points communs à $d$ et au cercle de centre $M$, rayon $3$.
\item Justifier géométriquement leur nombre avant de résoudre les équations.
\item Déterminer le lieu des centres des cercles passant par ces deux points.
\end{enumerate}

\section*{Problème S4 -- Contrôle et prédiction}

Une pièce est défectueuse avec probabilité $1/10$. Un test est positif
avec probabilité $4/5$ en cas de défaut et $1/5$ sinon. Les données sont fictives.
\begin{enumerate}
\item Calculer $P(+)$, $P_+(D)$ et $P_-(D)$.
\item Pour $X=\mathbf1_D$, donner l'espérance et la variance.
\item Si l'on connaît seulement le résultat du test, donner la meilleure
prédiction de $X$ au sens de l'erreur quadratique.
\item Vérifier que son erreur est plus faible que celle de la prédiction constante.
\end{enumerate}

\section*{Problème S5 -- Un jeu, une loi et des moyennes}

Quatre essais indépendants réussissent avec probabilité $1/4$.
Le nombre de réussites est $X$ ; le gain net est $G=5X-m$.
\begin{enumerate}
\item Donner la loi de $X$ et vérifier sa somme.
\item Déterminer la mise équitable et l'écart type du gain.
\item Pour $100$ parties indépendantes avec cette mise, donner l'écart type du gain moyen.
\item Majorer la probabilité que ce gain moyen s'écarte de $0$ d'au moins $1$.
\end{enumerate}

\section*{Problème S6 -- Approcher une racine et garantir l'erreur}

On veut approcher $\sqrt3$.
\begin{enumerate}
\item Donner un encadrement initial et une méthode de dichotomie.
\item Partir de $u_0=2$ avec Newton. Calculer $u_1,u_2$ exactement.
\item Construire un encadrement avec $u_2$ et sa largeur.
\item Comparer la preuve d'arrêt à une simple règle $|u_{n+1}-u_n|<\varepsilon$.
\end{enumerate}

\section*{Problème S7 -- Une droite comme prédicteur}

Pour les données $(0,1),(1,2),(2,2)$, on cherche un modèle affine.
\begin{enumerate}
\item Établir les équations normales par orthogonalité.
\item Résoudre et calculer l'erreur quadratique totale.
\item Comparer au meilleur modèle constant.
\item Prouver que la moyenne des valeurs prédites égale celle des observations.
\end{enumerate}

\section*{Problème S8 -- Un algorithme de moyenne et sa limite}

Deux quantités suivent
$x_{n+1}=(3x_n+y_n)/4$ et $y_{n+1}=(x_n+3y_n)/4$.
\begin{enumerate}
\item Étudier $s_n=x_n+y_n$ et $d_n=x_n-y_n$.
\item En déduire des formules explicites et les limites.
\item Interpréter la limite comme projection dans $\R^2$.
\item Pour $(x_0,y_0)=(10,2)$, déterminer un rang garantissant que les
deux composantes diffèrent de leur limite d'au plus $0{,}01$.
\end{enumerate}
